% ============================================================
\subsection{Markov Chains}
% ============================================================

A Markov chain models a system that moves between states in steps, where the rule for the next move depends only on where the system is now.

\begin{definition}[Markov Property]
A process has the Markov property when the distribution of the next state depends only on the current state. The path taken to reach the current state carries no additional information.
\end{definition}

This is memorylessness generalized from waiting times to states. The Geometric and Exponential forgot how long they had already waited. A Markov chain forgets its entire past trajectory and keeps only its current position, which is what makes the current state a sufficient summary for predicting the future.

\begin{definition}[States and Transitions]
A \textit{state} is a complete description of the system at one step. Each state carries \textit{transition probabilities} giving the chance of moving to each other state in one step. A \textit{transient} state is one the chain eventually leaves, and an \textit{absorbing} (terminal) state is one it can never leave, such as ``win'' or ``lose''. Collecting the transition probabilities into a matrix $P$, with entry $P_{ij}$ the probability of moving from state $i$ to state $j$, gives the \textit{transition matrix}. Each row sums to 1, since the chain must go somewhere.
\end{definition}

\vspace{1em}

\begin{keybox}[Two questions a Markov chain answers]
Problems split by whether the chain can get stuck.
\begin{itemize}[itemsep=1ex, label=--]
  \item \textbf{Absorbing chains} contain terminal states the process eventually lands in. The questions are which absorbing state it reaches, a \textit{hitting probability}, and how long it takes to get there, a \textit{hitting time}. Both are solved by first-step analysis.
  \item \textbf{Recurrent chains} keep moving with no absorbing state. The question is the long-run fraction of time spent in each state, given by the stationary distribution.
\end{itemize}
\end{keybox}

\subsubsection{Absorbing Chains and First-Step Analysis}

A problem with terminal states is solved by \textbf{first-step analysis}, which is the same conditioning-on-the-next-move idea behind the recursive expectation examples earlier. Condition on the single step the chain takes from its current state, and write the quantity at that state as a weighted average of the same quantity at the neighboring states, with the transition probabilities as weights. The recipe is:
\begin{enumerate}[itemsep=-1ex]
  \item enumerate the states and mark each as transient or absorbing
  \item for each transient state, write the quantity of interest as a weighted average over the states reachable in one step
  \item use the absorbing states to set the boundary conditions
  \item solve the resulting system of equations
\end{enumerate}
The quantity of interest is whichever the problem asks for. Filling it with a probability yields a hitting probability, and filling it with an expected step count yields a hitting time.

\newpage

\begin{example}[Gambler's Ruin]
A gambler starts with $k$ dollars and bets \$1 at a time, winning a dollar with probability $p$ and losing one with probability $q = 1-p$. Play continues until the stake reaches \$0 (ruin) or a target \$$N$. The state space is $\mathcal{S} = \{0, 1, \dots, N\}$, with $0$ and $N$ absorbing.

Let $u_i = P(\text{ruin} \mid \text{start at } i)$. Conditioning on the first bet gives the first-step equation
\[
u_i = p\, u_{i+1} + q\, u_{i-1},
\]
with boundary conditions set by the absorbing states,
\[
u_0 = 1 \quad (\text{already ruined}), \qquad u_N = 0 \quad (\text{already at target}).
\]
Writing $r = \frac{q}{p}$, the solution for a starting stake $k$ is
\[
u_k = \frac{r^k - r^N}{1 - r^N}.
\]
In the fair case $p = q = \tfrac{1}{2}$ this is undefined as written, and the equation instead solves to
\[
u_k = \frac{N - k}{N},
\]
so the probability of reaching the target is $\frac{k}{N}$, simply the stake as a fraction of the goal.
\end{example}

\vspace{1em}

\begin{keybox}[Expected time to absorption]
The expected number of steps until the chain is absorbed, starting from state $k$, is another first-step analysis. Writing $D_k$ for that expected duration,
\[
D_k = 1 + p\, D_{k+1} + q\, D_{k-1}, \qquad D_0 = D_N = 0,
\]
where the leading $1$ counts the current step. For the fair Gambler's Ruin ($p = q = \tfrac{1}{2}$) this solves to
\[
D_k = k(N - k).
\]
The product form matches intuition. Duration is longest when the gambler starts in the middle, far from both barriers, and shrinks to zero at either edge.
\end{keybox}

\subsubsection{Recurrent Chains and the Stationary Distribution}

When a chain has no absorbing states and keeps moving indefinitely, the relevant question is not where it ends but how its time is divided among the states over the long run.

\begin{definition}[Stationary Distribution]
For a chain with transition matrix $P$, a stationary distribution is a probability vector $\pi$ satisfying
\[
\pi P = \pi, \qquad \sum_i \pi_i = 1.
\]
It is the long-run fraction of time the chain spends in each state. The equation $\pi P = \pi$ says that a chain already distributed as $\pi$ reproduces that same distribution after one step.
\end{definition}

For a chain that is irreducible and aperiodic, meaning every state is reachable from every other and the chain does not cycle in a fixed period, this $\pi$ is unique and the chain settles into it regardless of where it started.